% Appendice aggiornata al 29/06/2026.
% Il documento principale deve caricare almeno il pacchetto listings:
% \usepackage{listings}

\section{Configurazione avanzata e report di contabilita}
\label{sec:configurazione-avanzata}

Questa sezione descrive le estensioni introdotte nel generatore di report e
spiega come configurare nuovi dataset e nuovi report senza modificare il
codice Python. Le funzionalita qui descritte sono compatibili con i file di
configurazione precedenti: una proprieta nuova viene applicata solamente
quando e presente nel file JSON.

\subsection{Riepilogo delle funzionalita aggiunte}

Le principali estensioni sono:

\begin{itemize}
  \item selezione esplicita della sorgente DBF ICR, CBR oppure locale;
  \item caricamento di un file di configurazione diverso da quello standard;
  \item filtri applicati direttamente durante la lettura dei DBF;
  \item filtri logici annidati mediante i gruppi \texttt{any} e \texttt{all};
  \item filtro di una sorgente in base ai valori gia letti da un'altra sorgente;
  \item trasformazione delle chiavi di join;
  \item copia di colonne gia esistenti;
  \item creazione di colonne derivate mediante prefissi, suffissi e
        selezione parziale del testo;
  \item creazione di colonne condizionali;
  \item sottrazione tra colonne numeriche;
  \item protezione dei testi che iniziano con il carattere \texttt{=},
        affinche Excel non li interpreti come formule;
  \item modalita di formattazione veloce per report con molte righe;
  \item generazione del report contabile con le colonne Dare, Avere e
        Differenza;
  \item distribuzione del programma tramite eseguibile e file batch.
\end{itemize}

\subsection{Avvio del programma}

Il programma si avvia dalla cartella principale del progetto con il comando:

\begin{lstlisting}[language=bash]
python .\script\main.py [opzioni]
\end{lstlisting}

Le opzioni disponibili sono le seguenti.

\begin{description}
  \item[\texttt{--config PERCORSO}]
  Indica il file JSON da utilizzare. Se non viene specificato, il programma
  usa \texttt{config/report\_config.json}.

  \item[\texttt{--report NOME}]
  Seleziona uno dei report presenti nell'oggetto \texttt{reports}. Se viene
  omesso, viene usato il valore di \texttt{default\_report}.

  \item[\texttt{--source icr}]
  Usa la sorgente ICR predefinita:
  \texttt{\textbackslash\textbackslash172.16.2.13\textbackslash arca\textbackslash ditte\textbackslash ICR}.

  \item[\texttt{--source cbr}]
  Usa la sorgente CBR predefinita:
  \texttt{\textbackslash\textbackslash172.16.2.13\textbackslash arca\textbackslash ditte\textbackslash CBR}.

  \item[\texttt{--source esempi}]
  Usa la cartella locale \texttt{esempi}. Questa opzione e utile per prove
  che non devono leggere la rete.

  \item[\texttt{--months N}]
  Indica quanti mesi elaborare per i report mensili. Il valore predefinito e
  3. L'opzione ha effetto sui dataset che calcolano il campo \texttt{Younth}.

  \item[\texttt{--no-open}]
  Genera il file Excel senza aprirlo automaticamente. Questa opzione e
  consigliata nei file batch e nelle esecuzioni pianificate.
\end{description}

Esempio di esecuzione del report contabile CBR:

\begin{lstlisting}[language=bash]
python .\script\main.py ^
  --config .\config\contabilita_config.json ^
  --source cbr ^
  --no-open
\end{lstlisting}

\subsubsection{Variabili di ambiente}

Due variabili di ambiente permettono di sostituire i percorsi definiti dal
programma:

\begin{description}
  \item[\texttt{PROGETTO\_CAVE\_DBF\_DIR}]
  Percorso della cartella che contiene i DBF. Se presente, ha precedenza
  sull'opzione \texttt{--source}.

  \item[\texttt{PROGETTO\_CAVE\_EXPORT\_DIR}]
  Cartella nella quale vengono salvati i file Excel.
\end{description}

Esempio PowerShell:

\begin{lstlisting}[language=bash]
$env:PROGETTO_CAVE_DBF_DIR = "\\172.16.2.13\arca\ditte\CBR"
$env:PROGETTO_CAVE_EXPORT_DIR = "C:\Report\Contabilita"
python .\script\main.py --config .\config\contabilita_config.json --no-open
\end{lstlisting}

\subsection{Struttura generale del file JSON}

Un file di configurazione contiene tre sezioni principali:

\begin{lstlisting}
{
  "default_report": "nome_report",
  "datasets": {
    "nome_dataset": {
      "...": "..."
    }
  },
  "reports": {
    "nome_report": {
      "dataset": "nome_dataset",
      "...": "..."
    }
  }
}
\end{lstlisting}

\begin{description}
  \item[\texttt{default\_report}]
  Nome del report eseguito quando non viene passato \texttt{--report}.

  \item[\texttt{datasets}]
  Descrive come leggere i DBF, collegare le tabelle e preparare le colonne.

  \item[\texttt{reports}]
  Descrive quali righe e colonne esportare e come costruire il file Excel.
\end{description}

Uno stesso dataset puo essere riutilizzato da piu report. In questo modo i
DBF e le join vengono definiti una sola volta, mentre ogni report puo avere
filtri, aggregazioni e colonne differenti.

\subsection{Definizione delle sorgenti DBF}

Le sorgenti sono dichiarate nell'oggetto \texttt{sources}:

\begin{lstlisting}
"sources": {
  "righe_contabili": {
    "file": "righe.DBF",
    "suffix": "",
    "columns": [
      "ID",
      "ID_TESTA",
      "CODCONTO",
      "DATAREG",
      "IMPORTOLIT"
    ]
  },
  "teste_contabili": {
    "file": "teste.DBF",
    "suffix": "teste",
    "columns": [
      "ID",
      "DESCRIZION"
    ]
  }
}
\end{lstlisting}

\begin{description}
  \item[\texttt{file}]
  Nome del DBF all'interno della cartella sorgente.

  \item[\texttt{columns}]
  Elenco dei campi da leggere. E consigliabile includere solamente i campi
  necessari, per ridurre memoria e tempo di elaborazione.

  \item[\texttt{suffix}]
  Suffisso aggiunto ai nomi dei campi per evitare collisioni tra tabelle.
\end{description}

Il nome prodotto da una sorgente segue la forma:

\begin{lstlisting}
NOME_CAMPO + "_" + suffix
\end{lstlisting}

Esempi:

\begin{itemize}
  \item \texttt{ID} con suffisso vuoto diventa \texttt{ID\_};
  \item \texttt{ID} con suffisso \texttt{teste} diventa
        \texttt{ID\_teste};
  \item \texttt{DESCRIZION} con suffisso \texttt{conti} diventa
        \texttt{DESCRIZION\_conti}.
\end{itemize}

Questi nomi completi devono essere usati nelle join e nella proprieta
\texttt{keep\_columns}.

\subsection{Filtri durante la lettura dei DBF}

I filtri di sorgente vengono applicati prima di creare il DataFrame. Sono
particolarmente importanti quando un DBF contiene centinaia di migliaia o
milioni di record.

\subsubsection{Filtro per data}

\begin{lstlisting}
"date_filter": {
  "column": "DATADOC"
}
\end{lstlisting}

Il valore del campo deve essere compreso tra la data iniziale e quella finale
del periodo elaborato.

\subsubsection{Filtro basato su un'altra sorgente}

\begin{lstlisting}
"filter_in": {
  "column": "ID_TESTA",
  "source": "doctes",
  "source_column": "ID"
}
\end{lstlisting}

Il record viene mantenuto solamente se \texttt{ID\_TESTA} appartiene
all'insieme dei valori \texttt{ID} gia letti dalla sorgente \texttt{doctes}.

\textbf{Importante:} le sorgenti sono lette nell'ordine in cui compaiono nel
JSON. La sorgente indicata da \texttt{source} deve quindi essere dichiarata
prima della sorgente che usa \texttt{filter\_in}.

Questa tecnica riduce notevolmente il numero di righe caricate. Nel report
contabile, ad esempio, vengono lette da \texttt{teste.DBF} solamente le teste
richiamate dalle righe contabili selezionate.

\subsubsection{Filtro logico \texttt{row\_filter}}

Il filtro \texttt{row\_filter} lavora direttamente sui valori del record DBF.
Puo contenere una condizione singola oppure gruppi annidati.

Il gruppo \texttt{all} equivale all'operatore logico SQL \texttt{AND}:

\begin{lstlisting}
"all": [
  {
    "column": "CODCONTO",
    "operator": "ge",
    "value": "02402905050"
  },
  {
    "column": "CODCONTO",
    "operator": "le",
    "value": "02402905570"
  }
]
\end{lstlisting}

Il gruppo \texttt{any} equivale all'operatore logico SQL \texttt{OR}:

\begin{lstlisting}
"any": [
  {
    "column": "CODCONTO",
    "operator": "starts_with",
    "value": "01301810"
  },
  {
    "column": "CODCONTO",
    "operator": "eq",
    "value": "01301830050"
  }
]
\end{lstlisting}

I gruppi possono essere annidati. Il seguente esempio rappresenta la
condizione SQL:

\begin{lstlisting}[language=SQL]
WHERE
    (
        CODCONTO >= '02402905050'
        AND CODCONTO <= '02402905570'
    )
    OR CODCONTO LIKE '01301810%'
    OR CODCONTO = '01301830050'
\end{lstlisting}

Configurazione equivalente:

\begin{lstlisting}
"row_filter": {
  "any": [
    {
      "all": [
        {
          "column": "CODCONTO",
          "operator": "ge",
          "value": "02402905050"
        },
        {
          "column": "CODCONTO",
          "operator": "le",
          "value": "02402905570"
        }
      ]
    },
    {
      "column": "CODCONTO",
      "operator": "starts_with",
      "value": "01301810"
    },
    {
      "column": "CODCONTO",
      "operator": "eq",
      "value": "01301830050"
    }
  ]
}
\end{lstlisting}

\subsubsection{Operatori disponibili}

Gli operatori supportati dai filtri sono:

\begin{description}
  \item[\texttt{eq}] uguale;
  \item[\texttt{ne}] diverso;
  \item[\texttt{ge}] maggiore o uguale;
  \item[\texttt{le}] minore o uguale;
  \item[\texttt{gt}] maggiore;
  \item[\texttt{lt}] minore;
  \item[\texttt{between}] compreso tra \texttt{min} e \texttt{max}, estremi
        inclusi;
  \item[\texttt{starts\_with}] il testo inizia con il valore indicato;
  \item[\texttt{in}] il valore appartiene all'elenco \texttt{values};
  \item[\texttt{blank}] il campo e vuoto;
  \item[\texttt{not\_blank}] il campo non e vuoto.
\end{description}

Esempio dell'operatore \texttt{between}:

\begin{lstlisting}
{
  "column": "CODCONTO",
  "operator": "between",
  "min": "02402905050",
  "max": "02402905570"
}
\end{lstlisting}

Esempio dell'operatore \texttt{in}:

\begin{lstlisting}
{
  "column": "CAUSALE",
  "operator": "in",
  "values": ["BA", "PF", "IF"]
}
\end{lstlisting}

\textbf{Attenzione ai codici con zeri iniziali:} nei filtri di sorgente i
valori vengono confrontati come testo ripulito dagli spazi. Codici conto,
codici articolo e causali devono quindi essere scritti tra virgolette.

\subsection{Sorgente principale e join}

La proprieta \texttt{base\_source} sceglie la tabella dalla quale parte il
dataset:

\begin{lstlisting}
"base_source": "schede_contabili"
\end{lstlisting}

Le altre sorgenti vengono collegate mediante \texttt{joins}:

\begin{lstlisting}
"joins": [
  {
    "source": "contabilita_teste",
    "left_on": "ID_TESTA_",
    "right_on": "ID_teste",
    "how": "inner"
  },
  {
    "source": "conti",
    "left_on": "CODCONTO_",
    "right_on": "CODICE_conti",
    "how": "left"
  }
]
\end{lstlisting}

\begin{description}
  \item[\texttt{source}]
  Nome della sorgente da collegare.

  \item[\texttt{left\_on}]
  Colonna presente nel DataFrame costruito fino a quel momento.

  \item[\texttt{right\_on}]
  Colonna della sorgente indicata da \texttt{source}.

  \item[\texttt{how}]
  Tipo di join Pandas. I valori normalmente usati sono \texttt{inner} e
  \texttt{left}. Se omesso, il valore predefinito e \texttt{left}.
\end{description}

\texttt{inner} elimina le righe prive di corrispondenza; \texttt{left}
conserva tutte le righe del lato sinistro e lascia vuoti i campi non trovati.

\subsubsection{Trasformazione delle chiavi di join}

Una chiave puo essere trasformata prima della join mediante
\texttt{left\_transform} o \texttt{right\_transform}.

\begin{lstlisting}
{
  "source": "conti_livello_1",
  "how": "left",
  "left_transform": {
    "column": "CODICE_conti",
    "left": 2,
    "suffix": "00000000"
  },
  "right_on": "CODICE_conti1"
}
\end{lstlisting}

Le trasformazioni disponibili sono:

\begin{description}
  \item[\texttt{left}] mantiene i primi $N$ caratteri;
  \item[\texttt{right}] mantiene gli ultimi $N$ caratteri;
  \item[\texttt{prefix}] aggiunge testo all'inizio;
  \item[\texttt{suffix}] aggiunge testo alla fine.
\end{description}

\subsection{Selezione e rinomina delle colonne}

\subsubsection{\texttt{keep\_columns}}

\texttt{keep\_columns} limita il dataset alle colonne necessarie:

\begin{lstlisting}
"keep_columns": [
  "ID_",
  "DATAREG_",
  "DAREFLAG_",
  "IMPORTOLIT_",
  "DESCRIZION_conti"
]
\end{lstlisting}

Questa operazione avviene dopo le join e prima delle rinomine. Una colonna non
presente viene ignorata; e comunque preferibile correggere il nome per non
nascondere errori di configurazione.

\subsubsection{\texttt{rename\_columns}}

\begin{lstlisting}
"rename_columns": {
  "DATAREG_": "Data Reg.",
  "DAREFLAG_": "Dare\\Avere",
  "IMPORTOLIT_": "importo_contabile",
  "DESCRIZION_conti": "Conto"
}
\end{lstlisting}

La chiave e il nome tecnico prima della rinomina; il valore e il nome usato
nelle fasi successive.

Nel JSON il carattere barra inversa deve essere raddoppiato. Per ottenere
\texttt{Dare\textbackslash Avere} si scrive quindi:

\begin{lstlisting}
"Dare\\Avere"
\end{lstlisting}

\subsection{Copia e derivazione delle colonne}

\subsubsection{Copia di una colonna}

\begin{lstlisting}
"copy_columns": {
  "id_schede_contabili": "id_contabilita_righe"
}
\end{lstlisting}

La chiave indica la colonna sorgente e il valore indica il nuovo nome. La
copia viene eseguita solamente se la sorgente esiste e la destinazione non
esiste gia.

\subsubsection{Derivazioni testuali}

\texttt{derived\_columns} permette di costruire nuove colonne senza scrivere
codice Python.

\begin{lstlisting}
"derived_columns": {
  "codice_conti1": {
    "source": "codice_conti",
    "left": 2,
    "suffix": "00000000"
  },
  "codice_conti2": {
    "source": "codice_conti",
    "left": 4,
    "suffix": "000000"
  }
}
\end{lstlisting}

Con \texttt{codice\_conti = "02402905350"} si ottiene:

\begin{itemize}
  \item \texttt{codice\_conti1 = "0200000000"};
  \item \texttt{codice\_conti2 = "0240000000"}.
\end{itemize}

L'ordine delle trasformazioni e: \texttt{left}, \texttt{right},
\texttt{prefix}, \texttt{suffix}. Se il valore sorgente e nullo, il risultato
e una stringa vuota.

\subsubsection{Colonna condizionale}

L'operazione \texttt{conditional} copia il valore della sorgente solamente
quando una condizione e verificata.

\begin{lstlisting}
"Dare": {
  "operation": "conditional",
  "source": "importo_contabile",
  "when": {
    "column": "Dare\\Avere",
    "equals": "D"
  },
  "otherwise": 0,
  "fillna": 0
}
\end{lstlisting}

Significato delle proprieta:

\begin{description}
  \item[\texttt{operation}]
  Deve valere \texttt{conditional}.

  \item[\texttt{source}]
  Colonna dalla quale copiare il valore.

  \item[\texttt{when.column}]
  Colonna da controllare.

  \item[\texttt{when.equals}]
  Valore atteso. Il confronto testuale ignora gli spazi iniziali e finali.

  \item[\texttt{otherwise}]
  Valore usato quando la condizione e falsa.

  \item[\texttt{fillna}]
  Valore usato quando il risultato e nullo. La proprieta e facoltativa.
\end{description}

Esempio per la colonna Avere:

\begin{lstlisting}
"Avere": {
  "operation": "conditional",
  "source": "importo_contabile",
  "when": {
    "column": "Dare\\Avere",
    "equals": "A"
  },
  "otherwise": 0,
  "fillna": 0
}
\end{lstlisting}

\subsubsection{Sottrazione tra colonne}

\begin{lstlisting}
"Differenza": {
  "operation": "subtract",
  "left": "Avere",
  "right": "Dare"
}
\end{lstlisting}

Il calcolo eseguito e:

\begin{lstlisting}
Differenza = Avere - Dare
\end{lstlisting}

Entrambe le colonne vengono convertite in numeri. Valori vuoti o non
convertibili vengono considerati zero.

\textbf{Ordine delle derivazioni:} le colonne derivate sono costruite
nell'ordine in cui sono dichiarate nel JSON. Nell'esempio contabile
\texttt{Dare} e \texttt{Avere} devono comparire prima di
\texttt{Differenza}.

\subsection{Campi calcolati predefiniti}

La proprieta \texttt{calculated\_fields} attiva i calcoli Python gia
implementati:

\begin{lstlisting}
"calculated_fields": ["Younth", "Cava", "Val_tot"]
\end{lstlisting}

\begin{description}
  \item[\texttt{Younth}]
  Converte la data documento nel periodo mensile.

  \item[\texttt{Cava}]
  Determina la cava dal prefisso del codice articolo e dal file
  \texttt{config/config\_cava.txt}.

  \item[\texttt{Val\_tot}]
  Calcola il valore totale della riga considerando prezzo, sconti, trasporto
  e quantita.
\end{description}

Un dataset che non necessita di questi calcoli usa:

\begin{lstlisting}
"calculated_fields": []
\end{lstlisting}

\subsection{Configurazione del report}

Ogni voce dell'oggetto \texttt{reports} punta a un dataset:

\begin{lstlisting}
"reports": {
  "schede_contabili": {
    "dataset": "schede_contabili",
    "output_file": "schede_contabili.xlsx",
    "main_sheet": "Schede"
  }
}
\end{lstlisting}

\subsubsection{Filtri finali}

\texttt{include\_rows} conserva le righe che soddisfano la condizione:

\begin{lstlisting}
"include_rows": [
  {
    "column": "codconto",
    "operator": "starts_with",
    "value": "01301810"
  }
]
\end{lstlisting}

\texttt{drop\_rows} elimina le righe che soddisfano la condizione:

\begin{lstlisting}
"drop_rows": [
  {
    "all": [
      {
        "column": "Codice_Articolo",
        "operator": "blank"
      },
      {
        "column": "Val_tot",
        "operator": "eq",
        "value": 0
      }
    ]
  }
]
\end{lstlisting}

Entrambe le proprieta supportano gli stessi gruppi \texttt{any}/\texttt{all}
e gli stessi operatori descritti per \texttt{row\_filter}.

La differenza e il momento di applicazione:

\begin{itemize}
  \item \texttt{row\_filter} agisce durante la lettura del DBF e riduce subito
        il volume dei dati;
  \item \texttt{include\_rows} e \texttt{drop\_rows} agiscono dopo join,
        rinomine e colonne derivate.
\end{itemize}

Se sono presenti piu elementi in \texttt{include\_rows}, vengono applicati in
sequenza: ogni elemento restringe ulteriormente il risultato. Lo stesso vale
per \texttt{drop\_rows}.

Nei filtri finali un valore JSON numerico, ad esempio \texttt{0}, attiva il
confronto numerico. Un valore tra virgolette, ad esempio \texttt{"0"}, viene
trattato come testo. Per codici con zeri iniziali si devono sempre usare le
virgolette.

\subsubsection{Raggruppamenti e aggregazioni}

Un report di dettaglio non raggruppato usa:

\begin{lstlisting}
"group_by": []
\end{lstlisting}

Un report aggregato puo usare:

\begin{lstlisting}
"group_by": [
  "Cava",
  "tipo_documento"
],
"aggregations": {
  "Quantita": "sum",
  "Val_tot": "sum",
  "Younth": "first"
}
\end{lstlisting}

Le aggregazioni sono quelle supportate da Pandas, ad esempio \texttt{sum},
\texttt{first}, \texttt{min}, \texttt{max}, \texttt{mean} e \texttt{count}.

\subsubsection{Colonne e ordinamento}

\texttt{columns} sceglie e ordina le colonne del foglio:

\begin{lstlisting}
"columns": [
  "Conto",
  "Data Reg.",
  "Data Comp.",
  "Nr. Riga",
  "Cod.Causali",
  "Causale",
  "Descrizione",
  "Note riga",
  "Dare",
  "Avere",
  "Differenza"
]
\end{lstlisting}

Una colonna indicata ma non presente viene ignorata. Durante lo sviluppo di
un nuovo config e quindi necessario controllare le intestazioni generate,
perche un errore di battitura non interrompe l'esecuzione.

L'ordinamento puo essere definito con:

\begin{lstlisting}
"sort_by": [
  "Data Reg.",
  "Nr. Riga"
]
\end{lstlisting}

Se \texttt{sort\_by} non e presente, il programma usa \texttt{group\_by}.

\subsection{Formattazione Excel}

Esempio:

\begin{lstlisting}
"output_file": "schede_contabili.xlsx",
"main_sheet": "Schede",
"freeze_panes": "A2",
"column_widths": {
  "Conto": 42,
  "Data Reg.": 14,
  "Dare": 18,
  "Avere": 18,
  "Differenza": 18
},
"wrap_columns": [
  "Descrizione",
  "Note riga"
],
"number_formats": {
  "Dare": "#,##0.00",
  "Avere": "#,##0.00",
  "Differenza": "#,##0.00"
}
\end{lstlisting}

\begin{description}
  \item[\texttt{output\_file}]
  Nome del file Excel. Nei report mensili puo contenere segnaposto, ad
  esempio \texttt{report\_\{periodo\_file\}.xlsx}.

  \item[\texttt{main\_sheet}]
  Nome del foglio principale.

  \item[\texttt{freeze\_panes}]
  Cella usata per bloccare righe e colonne.

  \item[\texttt{column\_widths}]
  Larghezza delle colonne.

  \item[\texttt{wrap\_columns}]
  Colonne nelle quali il testo deve andare a capo.

  \item[\texttt{number\_formats}]
  Formati numerici Excel applicati ai valori.
\end{description}

\subsubsection{Modalita veloce}

Per report molto grandi si puo aggiungere:

\begin{lstlisting}
"fast_format": true
\end{lstlisting}

La modalita veloce mantiene:

\begin{itemize}
  \item stile dell'intestazione;
  \item larghezza delle colonne;
  \item blocco dei riquadri;
  \item formati numerici dichiarati in \texttt{number\_formats}.
\end{itemize}

Per ridurre il tempo e la memoria non applica invece a ogni cella i bordi,
l'allineamento, le righe alternate e il ritorno a capo. E indicata per
esportazioni come la contabilita, che contiene oltre 150.000 righe.

\subsubsection{Protezione dalle formule involontarie}

Un testo letto da un DBF puo iniziare con il carattere \texttt{=}. Excel e
OpenPyXL potrebbero interpretarlo come una formula, anche quando si tratta di
una nota o di dati non strutturati.

Il writer forza ora questi valori al tipo testo. Ad esempio:

\begin{lstlisting}
=testo proveniente dal DBF
\end{lstlisting}

viene conservato integralmente e non produce piu il messaggio:

\begin{lstlisting}
Record rimossi: Formula dalla parte /xl/worksheets/sheet1.xml
\end{lstlisting}

La protezione si applica al foglio principale, ai fogli aggiuntivi e ai
riepiloghi.

\subsection{Riepiloghi e fogli aggiuntivi}

\subsubsection{Riepilogo}

\begin{lstlisting}
"summaries": [
  {
    "sheet_name": "Riepilogo",
    "group_by": [
      "Cava",
      "tipo_documento"
    ],
    "aggregations": {
      "Quantita": "sum",
      "Val_tot": "sum"
    },
    "number_formats": {
      "Quantita": "#,##0.00",
      "Val_tot": "#,##0.00"
    }
  }
]
\end{lstlisting}

Ogni elemento crea un nuovo foglio nello stesso file Excel.

\subsubsection{Foglio basato su un altro report}

\begin{lstlisting}
"extra_sheets": [
  {
    "report": "magalias",
    "main_sheet": "Alias"
  }
]
\end{lstlisting}

Il report indicato deve essere presente nello stesso oggetto
\texttt{reports}. Il programma ne carica il dataset e inserisce il risultato
come foglio aggiuntivo.

\subsection{Esempio contabile completo}

Il report contabile parte da \texttt{righe.DBF} e collega:

\begin{itemize}
  \item \texttt{teste.DBF} tramite \texttt{ID\_TESTA = ID};
  \item \texttt{causali.DBF} tramite \texttt{CAUSALE = CODICE};
  \item \texttt{anagrafe.DBF} tramite \texttt{CODICECF = CODICE};
  \item \texttt{CONTI.DBF} tramite \texttt{CODCONTO = CODICE}.
\end{itemize}

I campi principali letti da \texttt{righe.DBF} sono:

\begin{lstlisting}
"columns": [
  "ID",
  "ID_TESTA",
  "CODICECF",
  "CODCONTO",
  "NOTE",
  "CAUSALE",
  "DATAREG",
  "DATACOMP",
  "DAREFLAG",
  "NUMRIGA",
  "IMPORTOLIT"
]
\end{lstlisting}

\texttt{IMPORTOLIT} e l'importo nella valuta contabile di base.
\texttt{DAREFLAG} stabilisce il lato della registrazione:

\begin{itemize}
  \item \texttt{D}: l'importo viene scritto nella colonna Dare;
  \item \texttt{A}: l'importo viene scritto nella colonna Avere.
\end{itemize}

Le formule logiche sono:

\begin{lstlisting}
Dare       = IMPORTOLIT se DAREFLAG = "D", altrimenti 0
Avere      = IMPORTOLIT se DAREFLAG = "A", altrimenti 0
Differenza = Avere - Dare
\end{lstlisting}

Esempi:

\begin{center}
\begin{tabular}{lrrr}
\hline
\textbf{Flag e importo} & \textbf{Dare} & \textbf{Avere} &
\textbf{Differenza} \\
\hline
D; 329.786,51 & 329.786,51 & 0,00 & -329.786,51 \\
A; 33.729,59  & 0,00 & 33.729,59 & 33.729,59 \\
D; valore vuoto & 0,00 & 0,00 & 0,00 \\
\hline
\end{tabular}
\end{center}

La prima riga e stata verificata con la registrazione di apertura del
01/01/2026: nel DBF il valore e \texttt{IMPORTOLIT = 329786.51} e il flag e
\texttt{D}.

\subsection{Release unica e file batch}

La release non viene piu duplicata per ogni tipo di report. Il comando
\texttt{build\_release.bat} crea un eseguibile PyInstaller \emph{onefile}, che
viene collocato direttamente nella radice \texttt{Release}:

\begin{lstlisting}
Release\
|-- Crea_Report.exe
|-- Batch\
|   |-- _esegui_report.bat
|   |-- run_contabilita_CBR.bat
|   |-- run_report_cave_ICR.bat
|   |-- run_tutti_anagrafica_CBR.bat
|   `-- run_tutti_anagrafica_ICR.bat
|-- release_anagrafica_CBR\
|-- release_anagrafica_ICR\
|-- release_contabilita_CBR\
`-- release_report_cave_ICR\
\end{lstlisting}

Tutti i batch richiamano lo stesso \texttt{Crea\_Report.exe}. Ogni launcher
indica al batch comune il file JSON da caricare, la sorgente DBF e la cartella
di esportazione. Le configurazioni rimangono nelle cartelle delle rispettive
release. Per esempio, \texttt{run\_contabilita\_CBR.bat} usa la config in
\texttt{release\_contabilita\_CBR}, mentre
\texttt{run\_report\_cave\_ICR.bat} usa quella in
\texttt{release\_report\_cave\_ICR}.

Il batch comune:

\begin{enumerate}
  \item si posiziona nella propria cartella;
  \item crea la cartella dei log se non esiste;
  \item riceve dal launcher la sorgente DBF ICR oppure CBR;
  \item imposta la cartella di esportazione;
  \item avvia \texttt{Crea\_Report.exe} con il config richiesto;
  \item salva output ed errori in un file di log con data e ora;
  \item restituisce il codice di uscita dell'eseguibile.
\end{enumerate}

Per l'esecuzione pianificata si usa
\texttt{Batch\textbackslash run\_report\_periodico.bat}, che richiama il batch
totale \texttt{run\_tutti\_periodico.bat}. Il batch totale avvia tutti i report
anche quando uno dei precedenti termina con errore. Al termine, se sono
presenti errori, crea sul Desktop il file
\texttt{Crea\_Report\_ERRORI.txt} con i nomi dei report non completati e i
relativi codici di uscita. Se tutti i report riescono, elimina l'eventuale nota
rimasta da un'esecuzione precedente.

Per eseguire il report contabile e sufficiente avviare:

\begin{lstlisting}
Release\Batch\run_contabilita_CBR.bat
\end{lstlisting}

Un codice di uscita \texttt{0} indica che il report e stato completato.
In caso di errore, il primo controllo da fare e il file piu recente nella
cartella \texttt{logs} della release corrispondente.

\subsection{Procedura consigliata per creare un nuovo config}

\begin{enumerate}
  \item Identificare i DBF e i campi necessari.
  \item Scegliere la sorgente principale con \texttt{base\_source}.
  \item Dichiarare per prime le sorgenti usate da eventuali
        \texttt{filter\_in}.
  \item Applicare \texttt{date\_filter} e \texttt{row\_filter} il prima
        possibile per ridurre i dati letti.
  \item Definire le join usando i nomi completi di suffisso.
  \item Limitare il dataset mediante \texttt{keep\_columns}.
  \item Rinominare le colonne con \texttt{rename\_columns}.
  \item Aggiungere copie e derivazioni nell'ordine delle loro dipendenze.
  \item Configurare \texttt{include\_rows} e \texttt{drop\_rows} sui nomi
        successivi alla rinomina.
  \item Definire \texttt{group\_by} e \texttt{aggregations}, oppure usare un
        elenco vuoto per un report di dettaglio.
  \item Definire l'ordine finale con \texttt{columns}.
  \item Configurare larghezze e formati Excel.
  \item Attivare \texttt{fast\_format} per file molto grandi.
  \item Provare prima il comando con \texttt{--no-open}.
  \item Controllare numero di righe, intestazioni, formati e alcuni record
        noti prima di creare la release.
\end{enumerate}

\subsection{Errori comuni}

\begin{description}
  \item[DBF mancante]
  Controllare \texttt{file}, la sorgente scelta e
  \texttt{PROGETTO\_CAVE\_DBF\_DIR}.

  \item[Colonna assente nel file Excel]
  Controllare il nome originale, il suffisso, \texttt{keep\_columns},
  \texttt{rename\_columns} e infine \texttt{reports.columns}.

  \item[Join senza risultati]
  Controllare i tipi di dato, gli spazi, gli zeri iniziali e il tipo
  \texttt{inner}/\texttt{left}. Se necessario usare una trasformazione della
  chiave.

  \item[\texttt{filter\_in} non funzionante]
  Verificare che la sorgente richiamata sia dichiarata prima.

  \item[Condizione logica errata]
  Rappresentare sempre esplicitamente le parentesi SQL mediante gruppi
  \texttt{all} e \texttt{any}.

  \item[Importi nel lato sbagliato]
  Verificare i valori reali di \texttt{DAREFLAG}; per la contabilita CBR
  \texttt{D} significa Dare e \texttt{A} significa Avere.

  \item[Excel segnala formule rimosse]
  Assicurarsi di usare una release aggiornata. Il writer attuale salva come
  testo i valori DBF che iniziano con \texttt{=}.

  \item[Il file Excel e aperto]
  Se il file principale non puo essere sovrascritto, il programma crea una
  copia con data e ora nel nome.

  \item[Report molto lento]
  Ridurre i campi in \texttt{columns}, filtrare durante la lettura e attivare
  \texttt{fast\_format}.
\end{description}
